%!TeX encoding = UTF-8
%!TeX program = pdflatex
%!TeX spellcheck = it_IT

\documentclass[binding=0.6cm, noexaminfo]{sapthesis}\usepackage{microtype}
\usepackage[italian]{babel}
\usepackage[utf8]{inputenc}
\usepackage{hyperref}
\usepackage{listings}
\usepackage{xcolor} 
\usepackage{titlesec}
\titleformat{\chapter}[hang]
  {\normalfont\huge\bfseries}
  {}{0em}
  {}

\lstset{
    basicstyle=\ttfamily\small,
    breaklines=true,
    frame=single,
    numbers=left,
    numberstyle=\tiny\color{gray},
    keywordstyle=\color{blue},
    commentstyle=\color{green!50!black},
    stringstyle=\color{red},
    aboveskip=3em, 
    belowskip=3em, 
}

\title{Emulatore CUDA}
\author{Gabriele Guido}
\IDnumber{2144286}
\course{Laurea Triennale in Ingegneria Informatica}
\courseorganizer{Facoltà di Ingegneria dell'Informazione, Informatica e Statistica}
\AcademicYear{2023/2024}
\advisor{Prof. Giorgio Grisetti}
\authoremail{guido.2144286@studenti.uniroma1.it}
\copyyear{2026}
\thesistype{Tesi di Laurea}

\begin{document}

\frontmatter
\maketitle

\begin{abstract}
    La programmazione in CUDA \textit{(Compute Unified Device Architecture)} è il modo
    più efficace per scrivere programmi eseguibili su GPU \textit{(Graphics Processing Unit)}. L'architettura
    delle schede grafiche consente numerose operazioni parallele, riducendo i tempi di esecuzione e aumentando
    il \textbf{throughput} (attività completate nell'unità di tempo). Negli ultimi anni l'utilizzo
    di questa tecnologia si è diffuso ulteriormante
     grazie alla velocità di esecuzione di calcoli sulle matrici e programmare in questa 
    modalità è diventata quindi una competenza richiesta. Il suo 
    punto di forza tuttavia è ciò che rende la progettazione degli algoritmi meno semplice: i dati 
    sono elaborati in parallelo e senza ordine perciò durante l'esecuzione potrebbero verificarsi delle dipendenze 
    dati che, se non trattate correttamente, porterebbro a risultati sbagliati. In CUDA 
    abbandoniamo il concetto di esecuzione sequenziale: in ogni istante un numero finito di thread 
    esegue lo stesso codice su una porzione dell'input dato al programma.Senza un ordine tra questi 
    diventa necessario disegnare algoritmi in cui i thread lavorino in modo indipendente tra loro.  
    \\\\
    Il progetto nasce dal bisogno di programmare in CUDA anche senza una scheda grafica sul proprio PC e
    consiste in una librearia in cui ci sono le funzionalità essenziali per testare la correttezza di 
        un algoritmo destinato al calcolo parallelo su GPU. A livello 
        prestazionale non c'è alcun vantaggio, la vera utilità sta nella scrittura e verifica 
    del funzionamento di algoritmi destinati al calcolo parallelo.
    
\end{abstract}

\tableofcontents

\mainmatter
\chapter{CUDA}
Per capire come funziona l'esecuzione di un programma su una GPU dobbiamo abbandonare la 
concezione sequenziale e iniziare a pensare a cosa succede nello spazio più che nel tempo. Prenderemo 
come esempio di input da processare dei buffer di N numeri (interi). Supponiamo di voler eseguire 
una moltiplicazione per uno scalare S in ogni cella di \textit{dest[N]} e salvare 
il risultato su \textit{src[N]}. Su una CPU non serve altro
che iterare l'operazione su tutte le N posizioni del nostro vettore.
\begin{lstlisting}[language=C]
    for(int i = 0; i < N; i++)
        dest[i] += S * src[i];
\end{lstlisting}

Volendo ottenere lo stesso risultato su una GPU, iterare in questo modo non avrebbe senso: per sfruttare
la potenza di calcolo dobbiamo ottimizzare l'esecuzione. In questo esempio dobbiamo pensare che le 
operazioni di moltiplicazione stavolta vengono fatte su tutte le celle contemporaneamente: ogni thread 
esegue lo stesso codice su una singola posizione dell'array.
\begin{lstlisting}[language=C]
        dest[idx] += S * src[idx];
\end{lstlisting}
In questo modo il costo di questa semplice operazione passa da essere lineare \textit{(O(n))} a costante 
\textit{(O(1))}. 
\paragraph{Griglia, Blocchi e Thread}
Per poter gestire input con dimensione maggiore dei core (e quindi thread) che abbiamo a disposizione,
l'esecuzione viene spezzata in blocchi: ogni blocco è grande quanto il numero massimo di thread a disposizione.
L'input viene processato come una griglia di blocchi quindi CUDA mette a disposizione dei programmatori 
le informazioni come numero di thread (in un blocco), numero di blocco (nella griglia) e dimensione 
del blocco. In questo modo possiamo tenere conto dello spostamento spaziale della nostra esecuzione, processando
il nostro input come se fosse composto da tante "mattonelle" (da non confondere con le celle di un array).


\chapter{MyCuda}
La libreria MyCuda contiente delle funzionalità 
utili alla progettazione e simulazione del calcolo parallelo. L'esecuzione 
effettiva è gestita tramite un numero limitato di thread che simulano i core di una GPU.
Il funzionamento generale è gestito tramite una struttura dati \textit{MyCudaItem} al cui 
interno troviamo i buffer di memeoria (con rispettiva dimensione) e il puntatore alla funzione 
scritta in \textit{"pseudo-CUDA"} che verrà eseguta simulando i core
della GPU. Nella libreria la maggior parte delle funzioni sono 
delle operazioni ausiliare che facilitano la fase di sviluppo e verifica: gestione dei 2 buffer (creazione,
 stampa ed eliminazione), il calcolo di valori speciali 
 (dimensione griglia) e verifica di condizioni necessarie (dimensione buffer). 
Il cuore della libreria consiste nella funzione che simula la chiamata al \textit{Kernel} della GPU 
e tutta la gestione, tramite strutture dati 
dedicate, di indici, parametri e argomenti usati dai thread. 
Al di là quindi delle varie funzionalità operative, l'essenza di questa
libreria è la possibilità di scrivere funzioni (in C) similmente a come si farebbe usando CUDA. 
\\\\    
\textbf{Nota:} L'intera libreria è stata pensata per trattare array di interi ma 
la meccanica può essere facilmente estesa ad altri tipi di dato.    

\section{Parametri}
Nel header \textit{myCuda.h} troviamo le \textbf{macro} usate nella libreria. L'unica veramente rilevante è 
\texttt{BLOCK\_DIM}, che indica la dimensione di ogni blocco e quindi il numero massimo di thread generati.
Le altre sono destinate a debugging e test.    
\section{Strumenti principali}
\subparagraph{MyCudaItem}
Ho scelto di usare una struttura dati per salvare gli elementi più ricorrenti e
rendere tutte le funzioni consistenti con questo tipo di dato. Le funzioni di 
libreria avranno sempre, come parametro, un puntatore a un oggetto di questo tipo, lanciando un errore
in caso non sia rispettata questa sintassi.  
\begin{lstlisting}[language=C]
struct MyCudaItem{
    int buffer_size;
    int size;
    int* src_buffer;
    int* dest_buffer;
    KernelLaunchFn kernel_fn;
};
\end{lstlisting}

\subparagraph{MyCudaKernelLaunch}
In questa funzione vengono creati i thread che andranno a eseguire la funzione puntata dal campo di un oggetto
MyCudaItem. Per rispettare la struttura a griglia dell'esecuzione, i thread sono generati ad ogni iterazione 
di un singolo blocco, passandogli come funzione da eseguire proprio la funzione \textit{kernel\_fn}. Per gestire 
il passaggio dei parametri a ogni thread ho creato una struttura dati apposita,
\begin{lstlisting}[language=C]
struct ThreadArgs{
    int threadIdx;
    int blockIdx;
    void* dest;
    void* src;
    int size;
};  
\end{lstlisting}
inizializzata con \texttt{MyCudaSetThreadArgs()}. In questo modo ogni thread potrà accedere alla sua 
posizione e ai buffer. Con \texttt{MyCudaIterations()} vengono calcolati i passaggi totali per consumare 
l'input, ovvero il numero di blocchi totali.  


\begin{lstlisting}[language=C]
void MyCudaKernelLaunch(MyCudaItem* cuda_item){
    assert(cuda_item && "cuda_item not valid");
    int passes = MyCudaIterations(cuda_item);
    pthread_t threads[BLOCK_DIM];
    ThreadArgs thread_arguments[BLOCK_DIM];
    for(int bid = 0; bid<passes;bid++){
            for(int tid = 0; tid < BLOCK_DIM; tid++){
                MyCudaSetThreadArgs(&thread_arguments[tid],tid,bid,cuda_item->dest_buffer,cuda_item->src_buffer,cuda_item->size);
                //kernel launch:
                pthread_create(&threads[tid],NULL,cuda_item->kernel_fn,&thread_arguments[tid]);
            }
            for(int tid = 0; tid < BLOCK_DIM; tid++){
                pthread_join(threads[tid],NULL);
            } 
        }
}
\end{lstlisting}
Una volta entrati nella funzione, i parametri vengono estratti dalla \texttt{struct} e salvati localmente.
Per vedere il loro utilizzo prendiamo l'esempio di \texttt{ReductionSum}, in cui sommiamo tutti i valori 
di un vettore. Nella sua versione \texttt{Kernel} troviamo
\begin{lstlisting}[language=C]
void* MyCudaReductionKernel(void* args){
    ThreadArgs* info = (ThreadArgs*) args;

    //estraggo i parametri 
    int idx = info->threadIdx+info->blockIdx*BLOCK_DIM;
    int size = info->size;
    int* src = info->src;
    int* dest = info->dest;

    //verifico se sono ancora nel buffer
    if(idx >= size){
        pthread_exit(NULL);
    }

    //sommo due celle vicine 
    dest[idx] = src[2*idx]+src[2*idx+1];
    pthread_exit(NULL);
}
\end{lstlisting}

\subparagraph{MyCudaMemSwap}
In molti algoritmi, prima di arrivare al risultato, ci sono diversi passaggi intermedi in cui si salvano 
risultati da usare successivamente. Per evitare di consumare molto spazio usiamo i due array di sorgente e 
destinazione in modo alternato. Questa funzione non fa altro che invertire, all'interno della struttura 
MyCudaItem, i puntatori alle due zone di memoria. Ad esempio nel caso di \texttt{ReductionSum} gli array 
vengono alternati finchè la dimensione finale è 1 (la somma totale). 

\begin{lstlisting}[language=C]

void MyCudaMemSwap(MyCudaItem* cuda_item){
    assert(cuda_item && "CudaItem not valid");
    int* aux = cuda_item->src_buffer;
    cuda_item->src_buffer = cuda_item->dest_buffer;
    cuda_item->dest_buffer = aux;
    for(int i = 0; i<cuda_item->buffer_size;i++){
        cuda_item->dest_buffer[i] = 0; 
    }
}
\end{lstlisting}




\backmatter
% Qui inserisci la bibliografia

\end{document}